\section{Selective attention under competing distractors}
\label{sec:distractor}

The base task contains no competing events. To ask how the coupled architecture
behaves when it must \emph{select} a target against a salient irrelevant event ---
the defining demand of natural selective attention
\citep{desimone1995neural, sprague2018rescuing, cosman2018prefrontal} --- we trained
the same shared-memory cross-talk model (internal identifier \texttt{v11\_part5},
$1{,}015{,}656$ parameters) on the distractor variant of
Section~\ref{sec:env-distractor}, in which the uncued side may undergo its own
task-irrelevant change that the agent must ignore. The analysis reveals two
selectivities riding on the single shared memory --- one for the \emph{decision},
one for \emph{valuation} --- and yields the clearest illustration in this work of the
$\mu$/$q$ dissociation.

\subsection{A value gate over a near-perfect detector}
Aggregate performance understates the agent (accuracy $\approx 0.60$, hit
$\approx 0.54$; Figure~\TODO{fig:dist-behavior}) because the aggregate averages over a
deliberate refusal to engage on low-value trials. Split by cue colour --- which sets
reward magnitude --- the picture is sharp: on high-value trials the agent is a
near-perfect detector (hit $\approx 0.83$, false-alarm rate $<0.01$,
$d'\approx 3.4$--$3.9$), whereas on the lowest-value trials it \emph{abandons the
task entirely}, declaring almost never (hit $\approx 0.00$) and bailing at cue onset.
The collapse is therefore not a perceptual failure but a value-driven
\emph{engagement} decision: when the reward on offer is too small, the agent declines
to work for it. Tellingly, the same architecture trained on the distractor-\emph{free}
task does \emph{not} abandon low-value trials --- the bail is specific to the
cost/benefit of selecting under competition. This is the model's analogue of
value-dependent task engagement, the kind of cost-of-control computation attributed
to anterior cingulate and ventral striatal circuits
\citep{maunsell2004neuronal, stanisor2013neural}.

\subsection{Two selectivities for one decision: a causal dissociation}
Clamping the attention of one pathway at the change frame and reading off behaviour
and value separately produces a double dissociation
(Figure~\TODO{fig:dist-causal}). Biasing the \emph{$\mu$ (policy) pathway} toward the
cued location swings the decision across its full range (hit rate $0.01 \to 1.00$)
while leaving the critic's value estimate essentially unchanged; biasing the \emph{$q$
(value) pathway} moves the value estimate but never the decision. The two pathways
share only the deep memory $\Htwo$, yet one commands the decision and the other the
valuation --- a within-model demonstration that the split readout carries genuine
functional content, not nominal labels.

\subsection{The distractor is filtered in the decision pathway, not globally}
Where is the irrelevant event suppressed? Not at the input: the distractor's location
is highly decodable from the fast sensory memory $\Hone$ (decoding accuracy
$\approx 0.94$), so the agent plainly \emph{sees} it. It is suppressed specifically in
the \emph{$\mu$ (decision) pathway}, whose attention to the distractor is near zero
(priority weight $\approx 0.004$ versus $\approx 0.26$ for the cued primary), while the
$q$ (value) pathway continues to read it. The suppression is causal: forcing the $\mu$
pathway's attention onto the distractor converts the agent from a competent detector
into one that false-alarms to the distractor (hit $1.00 \to 0.01$, false-alarm
$0.00 \to 0.99$; $d'$ inverts from $+5.5$ to $-4.9$). The distractor is thus
``seen but not acted upon'': filtered in the pathway that issues the response rather
than removed from the representation. This is a precise model counterpart of the
superior-colliculus ``dual effect'' --- unilateral SC inactivation simultaneously
impairs responding to relevant events and \emph{releases} responding to competing
ones --- in which the colliculus governs which signals inform the decision rather than
which are sensed \citep{lovejoy2010inactivation, zenon2012attention, herman2020attention}.

\subsection{Value reaches into the decision filter, but as an engage/abandon switch}
Finally, value does not merely set a post-perceptual response criterion: it reshapes
the decision filter itself. On high-value trials the $\mu$ pathway's allocation is
sharp and the distractor strongly suppressed; on low-value trials the allocation is
diffuse and leaky (cued/distractor weights $\approx 0.26/0.004$ for high value versus
$\approx 0.03/0.11$ for low value). Crucially, this modulation is \emph{binary} rather
than graded over reward magnitude --- the two engaged value levels behave alike, and
the difference is whether the agent engages at all. Value, in this model, gates
\emph{whether} the priority filter is deployed, not how finely it is set, correcting a
naive reading in which reward would simply scale a continuous criterion. The critic's
value estimates are themselves well calibrated and ordered by reward
(calibration slope $\approx 1.0$), consistent with a value system that prices the trial
and a decision system that, gated by that price, either deploys selective attention or
declines to.

\subsection{Summary and predictions}
The distractor variant shows the coupled architecture doing genuine selection: a
near-perfect detector whose decisions are commanded by an image-grounded $\mu$ pathway
that suppresses competitors at the level of action, gated by a $q$ value pathway that
prices the trial and decides whether to engage. The mechanism makes two further
predictions for the brain. First, a competing distractor should be \emph{more}
decodable from value/valuation populations than from the choice-driving populations
that suppress it --- the inverse of the usual assumption that suppression is global.
Second, because suppression lives in the decision pathway, perturbing the priority
node (SC/FEF) should release distractor responding while sparing the distractor's
sensory representation, whereas perturbing the value circuit should alter engagement
and value-modulated filtering without abolishing detection
\citep{zenon2012attention, herman2018gain, gregoriou2012lesions}.
